\section{GitHub}

\begin{frame}[t]
	\frametitle{\insertsection}

	\begin{columns}
		\begin{column}{.4\paperwidth}
			\begin{block}{Ventajas de GitHub Classroom}
				\begin{itemize}
					\item \textbf{Repositorio público} para cada estudiante.
					\item \textbf{Control de versiones} con Git.
					\item Tendrán registro completo de su trabajo.
					\item Facilita la evaluación y retroalimentación.
				\end{itemize}
			\end{block}
		\end{column}
		\begin{column}{.4\paperwidth}
			\begin{figure}[ht!]
				\centering
				\includegraphics[width=0.25\paperwidth]{githubclassroom}
				\caption{Interfaz de GitHub Classroom (Fuente: \url{https://classroom.github.com})}
				\label{fig:githubclassroom}
			\end{figure}
		\end{column}
	\end{columns}

	\begin{columns}
		\begin{column}{.4\paperwidth}
			\begin{block}{GitHub Desktop}
				\begin{itemize}
					\item Aplicación de escritorio que facilita el uso de Git y GitHub
					      sin necesidad de usar la línea de comandos.
					\item Aprenderán herramientas profesionales de desarrollo.
				\end{itemize}
			\end{block}
		\end{column}
		\begin{column}{.4\paperwidth}
			\begin{figure}[ht!]
				\centering
				\includegraphics[width=0.25\paperwidth]{githubdesktop}
				\caption{Interfaz de GitHub Desktop (Fuente: \url{https://desktop.github.com})}
				\label{fig:githubdesktop}
			\end{figure}
		\end{column}
	\end{columns}
\end{frame}
